\begin{table}[h]
    \centering
    \small
    \resizebox{\linewidth}{!}{
    \begin{tabular}{c|c|ccc}
    % \hline
    \toprule
         \multirow{3}{*}{Model}
        & \multirow{3}{*}{\#Params} 
        & \multirow{3}{*}{\shortstack{Statistical \\ Learning }} & \multirow{3}{*}{Fusion} & \multirow{3}{*}{\shortstack{Matching \\ Score \\ as features}} \\
         ~ & & & \\
         ~ & & & \\
         \hline
         SSE-Cross  & — & \xmark & \cmark & \xmark \\
         D\&R Net  & — & \cmark & \cmark & \xmark \\
         MCR & 2.33M & \xmark & \cmark & \xmark \\
         SANCL & 2.38M & \cmark & \cmark & \xmark \\
         CMCR & 2.41M & \xmark & \cmark & \xmark \\
         GBDT & 21.8M & \xmark & \cmark & \xmark \\
         \modelname  & 2.28M &  \xmark & \xmark  & \cmark \\
        \bottomrule
    \end{tabular}
    }
    \caption{Comparison between baseline models and ours.}
    \label{tab:model_comp}
\end{table}
\renewcommand{\arraystretch}{1.1}
\begin{table*}[ht]
    % \centering
    \small
    \resizebox{\textwidth}{!}{
    \begin{tabular}{l*{3}{|c c c}}
        \toprule
        \multirow{2}{*}{Model} & 
        \multicolumn{3}{c|}{Cloth. \& Jew.} & \multicolumn{3}{c|}{Electronics}  & \multicolumn{3}{c}{Home \& Kitchen}  \\
        % \cline{3-11}
         ~ & \textbf{MAP} & \textbf{N@3} & \textbf{N@5} & \textbf{MAP} & \textbf{N@3} & \textbf{N@5} & \textbf{MAP} & \textbf{N@3} & \textbf{N@5} \\
        \midrule
        SSE-Cross & 65.0 & 56.0 & 59.1 & 53.7 & 43.8 & 47.2 & 60.8 & 51.0 & 54.0 \\
        D\&R Net &  65.2 & 56.1 & 59.2 & 53.9 & 44.2 & 47.5 & 61.2 & 51.8 & 54.6 \\
        MCR  & 66.4 & 57.3 & 60.2 & 54.4 & 45.0 & 48.1 & 62.6 & 53.5 & 56.6  \\
        SANCL & 67.3 & 58.6 & 61.5 & 56.2 & 47.0 & 49.9 & 63.4 & 54.3 & 57.4 \\
        CMCR & 67.4 & 58.6 & 61.6 & 56.5 & 47.6 & 50.8 & 63.5 & 54.6 & 57.8 \\
        GBDT & 82.6 & 80.3 & 79.3 & 74.2 & 68.0 & 69.8 & 81.7 & 76.5 & 78.8 \\
        \modelname~(Ours) & \textbf{92.3} & \textbf{90.4} & \textbf{91.5} & \textbf{81.4} & \textbf{78.6} & \textbf{75.6} & \textbf{88.6} & \textbf{88.3} & \textbf{88.4}  \\
        \hline
        MCR+BERT & 65.8 & 55.9 & 58.8 & 55.9 & 46.8 & 49.4 & 62.4 & 52.9 & 56.1 \\
        GBDT+BERT & 80.3 & 78.7 & 77.1 & 73.8 & 68.3 & 69.5 & 81.4 & 76.9 & 79.4\\
        \modelname~+BERT~(Ours) & \textbf{91.5} & \textbf{89.7} & \textbf{90.1} & \textbf{79.3} & \textbf{77.7}& \textbf{78.6} & \textbf{87.7} & \textbf{85.9} & \textbf{86.2} \\
        \bottomrule
    \end{tabular}
    }
    \caption{Results on the Amazon-MRHP (English) dataset. All reported metrics are the average of five runs. Baseline results are from \citet{nguyen2023gradient}. PREMISE outperform the strongest baseline with p-value<0.05 based on the paired t-test.}
    \label{amazon results}
\end{table*}

\renewcommand{\arraystretch}{1.05}
\begin{table*}[ht]
    \centering
    \small
    \resizebox{\linewidth}{!}{
    \begin{tabular}{l*{3}{|c c c}}
        \toprule
         \multirow{2}{*}{Model} &
         \multicolumn{3}{c|}{Cloth. \& Jew} & \multicolumn{3}{c|}{Electronics}  & \multicolumn{3}{c}{Home \& Kitchen}   \\
        ~ & \textbf{MAP} & \textbf{N@3} & \textbf{N@5} & \textbf{MAP} & \textbf{N@3} & \textbf{N@5} & \textbf{MAP} & \textbf{N@3} & \textbf{N@5} \\
        \midrule
        SSE-Cross & 66.1 & 59.7 & 64.8 & 76.0 & 68.9 & 73.8 & 72.2 & 66.0 & 71.0 \\
        D\&R Net &  66.6 & 61.3 & 65.8 & 76.5 & 69.5 & 74.3 & 72.7 & 66.7 & 71.8 \\
        MCR &  68.8 & 62.3 & 67.0 & 76.8 & 70.7 & 75.0 & 73.8 & 67.0 & 72.2  \\
        SANCL & 70.2 & 64.6 & 68.8 & 77.8 & 71.5 & 76.1 & 75.1 & 68.4 & 73.3 \\
        CMCR & 70.3 & 64.7 & 69.0 & 78.2 & 72.4 & 79.6 & 75.2 & 68.8 & 73.7 \\
        GBDT & 78.5 & 77.1 & 79.0 & 87.9 & 86.7 & 88.1 & 85.6 & 78.8 & 83.1 \\
        \modelname~(Ours) &  \textbf{95.3} & \textbf{94.6} & \textbf{95.0}  & \textbf{96.9} & \textbf{96.1} & \textbf{96.8} & \textbf{93.9} & \textbf{91.5} & \textbf{92.8} \\
        % \cline{2-11}
        % & MCR+BERT &   \\
        % & \modelname~+BERT~(Ours) &  \\
        \bottomrule
    \end{tabular}
    }
    \caption{Results on the Lazada-MRHP (Indonesian) dataset. 
    % Notations have the same meaning as in Table \ref{amazon results}.
    }
    \label{lazada results}
\end{table*}


\section{Experiments}
\subsection{Datasets and Metrics}
We evaluate our model on two \taskabbrname~datasets \citep{liu2021multi}, each of which subsumes same three categories: \textit{Clothing, Shoes \& Jewelry, Home \& Kitchen} and \textit{Electronics}. We train and test both datasets on a single NVIDIA RTX A6000 GPU.
The gradients are calculated and backpropagated for each batch in a single forward pass, without batch division and gradient accumulation.
We compare our model with several baseline models on three common metrics for ranking tasks: the mean average precision (MAP), the N-term ($N=3,5$ in our experiment in accord with previous works) Normalized Discounted Cumulative Gain (NDCG@N)~\citep{jarvelin2017ir, diaz2018modeling}.
The helpfulness scores are labeled as the logarithm of approval votes to the corresponding reviews and are clipped into integers within $[0,4]$. 
The statistics of datasets and more training details are provided in appendix.


\subsection{Baselines}
We compare our model with the following baselines: Stochastic Shared Embeddings enhanced cross-modal network (SSE-Cross)~\cite{abavisani2020multimodal}, Decomposition and Relation Network~(D\&R Net)~\cite{xu2020reasoning}. The Multimodal Coherence Reasoning network~(MCR)~\cite{liu2021multi} designs several reasoning modules based on fused representations for prediction. SANCL~\cite{han-etal-2022-sancl} and contrastive-MCR (CMCR)~\cite{nguyen-etal-2022-adaptive} minimize auxiliary contrastive loss to refine the multimodal representations. 
Gradient-boosted decision tree (GBDT)~\cite{nguyen2023gradient} design a random walk policy and aggregate the helpfulness scores through from tree leaves---the endpoints of the random walk.

To provide a holistic view on the distinctions between the learning paradigms of these models, we list and compare three key characteristics between~\modelname~and baselines in Table~\ref{tab:model_comp}.
From the table we observe that all baseline models contain fusion modules inside the entire structures.
Moreover, D\&R Net and SANCL also incorporate extra statistical correlations (Adjective-Noun Pairs and selective-attention mask creation) that inject external knowledge to bridge the semantic gap between textual and visual modality or cast more focus on contents that perceived important by human beings.
Our model escapes from both complicated manually crafted features and conventional model architecture by directly computing the matching scores and automatically picking the $K$ highest ones as features for regression. 

\subsection{Results}
We run our models three times and report the average performance in Table~\ref{amazon results}~and~\ref{lazada results}.
It can be clearly seen that our model outperforms all these baselines on two datasets.
Particularly, compared with the strongest baseline---GBDT~\cite{nguyen2023gradient}, PREMISE gains over 5 points improvement on MAP and NDCG@5 and 10 points improvement on NDCG@3 on Amazon-MRHP dataset, and 6.8$\sim$17.5 improvement on all metrics on Lazada-MRHP datasets.
When using BERT to initialize embeddings, we note a slight performance degradation compared to the implementations that use GloVe as embeddings in both~\modelname and other baselines. 
Such outcome demonstrates the superiority of our fusion-free model and, at least in the MRHP task, multimodal fusion is not a necessity and may hinder the model from better performance.

Besides, we highlight the size of the feature vectors in the last layer.
Fusion-based baselines usually concatenate representations from both fields and modalities to perform final regression, which requires at least 512 (128$\times$4) dimensions of feature vectors.
The vector is even longer in MCR and SANCL (over 1000) since there are many extra features taken into account.
Nevertheless, the feature vector lengths for the best performance in~\modelname~are apparently smaller.
As shown in~\Cref{fig:perf_kvalue}, the optimal choices of $K$ range from 64 to 128.
This fact shows that there could be many redundant elements in the vectors generated by fusion-based models, and~\modelname~successfully enhance the efficiency of unit length features through a simple representation learning policy.

\begin{table}[ht]
    \small
    \centering
    \resizebox{\linewidth}{!}{
    \begin{tabular}{l|*{3}{c}}
        \toprule
        Description & \textbf{MAP} & \textbf{N@3} & \textbf{N@5} 
        \\
        \midrule
        \modelname~(Amazon) &  \textbf{87.4} & \textbf{85.8} & \textbf{86.5} 
        \\
        \midrule
            -w/o n-gram token repr & 84.8 & 82.1 & 83.3 \\
            -w/o sent repr & 86.2 & 84.3 & 85.0  \\
            -w/o n-gram sent repr & 85.7 & 83.9 & 84.6 \\
            -w/o n-gram RoI repr & 83.9 & 81.8 & 82.6 \\
            -w/o image repr & 86.5 & 84.1 & 85.3 \\
            -w/o n-gram token \& n-gram RoI repr & 75.3 & 69.8 & 72.2 \\
            -w/o n-gram sent repr \& image repr & 84.1 & 82.5 & 83.0 \\
        \midrule
        \modelname~(Lazada) &  \textbf{95.4} & \textbf{94.0} & \textbf{94.9} 
        \\
        \midrule
            -w/o n-gram token repr & 91.0 & 88.9 & 89.5 \\
            -w/o sent repr & 93.5 & 91.6 & 92.2 \\
            -w/o n-gram sent repr & 94.3 & 92.9 & 93.8 \\
            -w/o n-gram RoI repr & 92.7 & 90.1 & 91.7 \\
            -w/o image repr & 94.8 & 94.2 & 94.6 \\
            -w/o n-gram token \& n-gram RoI repr & 80.1 & 78.6 & 79.5  \\
            -w/o n-gram sent repr \& image repr & 92.3 & 89.6 & 90.1 \\
        \bottomrule
    \end{tabular}
    }
    \caption{Ablation experiments of~\modelname~on two datasets. The values are averaged over all three categories.}
    \label{tab:abl}
\end{table}

\subsection{Ablation Study}
We run our models several ablative settings under feature selection.
These settings correspond to excluding some features from different scales when computing the multiscale matching scores;
During implementation, we mask out these vectors. 

The results are summarized in Table~\ref{tab:abl}, from which we have the following discoveries.
First, discarding representations of any scale causes degradation in the model's performance, indicating that all these chosen features contribute to the accurate prediction.
In addition, performance plummets more severely when features are removed on smaller scales (or in bottom layers), including a single scale (e.g. ``n-gram token repr" vs. ``n-gram sent repr", ``n-gram RoI repr" v.s. ``n-gram image repr") or the combinations (e.g., ``n-gram token \& n-gram RoI repr" v.s. ``n-gram sent repr \& image repr``). 
This outcome reveals that lower-level feature is more fundamental to the model's performance than higher ones, since a large portion of matching scores are computed from them.
% such as the n-gram token representation or RoI representation.



